\Needspace{34\baselineskip}
\subsubsection*{Empirical Poisson laws for almost every realization}
\phantomsection\label{sec:empirical-poisson}
\pdfbookmark[2]{Empirical Poisson laws for almost every realization}{bookmark:empirical-poisson}
The comparison above averages over the prime signs. A summable error
also permits a different order of quantifiers: first sample the entire
function $f$, then fix it and randomize only the spatial origin. The
following consequence concerns a count, not a comparison of the whole
local field after $f$ has been fixed.

\begin{empiricalcorollary}[Spatial empirical law for a fixed realization]
\label{cor:empirical-poisson}
Fix $0<\alpha<1$ and $\tau>0$. With $V_M$ as in
\eqref{eq:saddle-definitions}, for all sufficiently large $k$ set
\[
 M_k=2^k,\qquad
 L_k=k-\left\lfloor\frac{\alpha V_{M_k}}{\log2}\right\rfloor,
 \qquad p_k=2^{-L_k},
\]
\[
 n_k=M_k-L_k,\qquad h_k=\lfloor\tau 2^{L_k}\rfloor,
 \qquad N_k=n_k-h_k+1.
\]
Let $U_k$ be uniform on $\{0,\ldots,n_k-h_k\}$, independently of $f$, and put
\[
 Z_k(f,u)=\sum_{j=u+1}^{u+h_k}J_{j+1,L_k}(f),\qquad
 \mu_k^f=\frac1{N_k}\sum_{u=0}^{n_k-h_k}\delta_{Z_k(f,u)}.
\]
Then, for almost every $f$ in the model \eqref{eq:model},
\[
 \dTV\!\left(\mu_k^f,\Pois(\tau)\right)
 =\dTV\!\left(\Law_{U_k}(Z_k(f,U_k)),\Pois(\tau)\right)
 \longrightarrow0.
\]
Here $h_k$ counts potential start sites, the two signs are aggregated,
and a run is counted only at its left-maximal start. In particular,
\[
 h_kp_k\longrightarrow\tau,\qquad
 h_k=M_k\exp\{-\alpha V_{M_k}+O(1)\}=M_k^{1-o(1)}.
\]
The full-probability set is asserted for the fixed pair $(\alpha,\tau)$;
the conclusion is along the specified dyadic scales.
\end{empiricalcorollary}

The starts read by $Z_k$ range from $u+2$ to $u+h_k+1$. Thus every
base of length $L_k$ is contained in the global prefix $[1,M_k]$, and
the initial run from $1$ is excluded. A full run is not censored at
$M_k$, although its excess is not used by this count. In
$\Law_{U_k}$ only the origin is random; no independence between the
origins at different scales is required.

\begin{proof}
We first check the arithmetic budget and then do the frequency
calculation solely in the independent target. The hard-scale
expansions in \cref{prop:cutoffs} give
\[
 V_{M_k}\asymp\sqrt{k\log k},\qquad
 \nu_{M_k}\asymp\sqrt{k/\log k},\qquad
 V_{M_k}/\nu_{M_k}\longrightarrow\infty.
\]
Writing $\Lambda_k=n_kp_k$, the definitions imply
\[
 \log\Lambda_k=\alpha V_{M_k}+O(1),\qquad
 h_k/n_k=\exp\{-\alpha V_{M_k}+O(1)\}.
\]
Hence $\Lambda_k\to\infty$, $1\le h_k\le n_k$ eventually, and
\eqref{eq:micro-budget} holds with $\mathcal A_{M_k}$ equal to the whole
probability space and any fixed $c>0$, for all sufficiently large $k$.
\Cref{thm:micro-run-tv} consequently gives a sequence $\Delta_k$ with
\[
 \dTV\bigl(\Law(\mathbf I_{M_k}),\Law(\mathbf P_{M_k})\bigr)
 \le\Delta_k,\qquad
 \sum_k\Delta_k<\infty,\qquad
 \sum_k h_k/n_k<\infty.
\]
Indeed, for fixed admissible $c'$ and $\varepsilon$, the bound is
$O(e^{-c'\nu_{M_k}}+2^{-k(1/3-\varepsilon)})$, and both terms are
summable.

Sum the target over the signs and excesses at a site. Since
$\sum_{e,\sigma}\gamma_{e,\sigma}=1$, the resulting variables
\[
 X_{k,j}=\sum_{e\ge0,\,\sigma=\pm1}P_{k,j,e,\sigma}
 \sim\Pois(p_k),\qquad 1\le j\le n_k,
\]
are independent. For $0\le u<N_k$ and $r\in\NN$, define
\[
 S_{k,u}=\sum_{j=u+1}^{u+h_k}X_{k,j},\qquad
 F_{k,r}=\frac1{N_k}\sum_{u=0}^{N_k-1}\1_{\{S_{k,u}=r\}},
 \qquad q_{k,r}=\Pois(h_kp_k)\{r\}.
\]
Every $S_{k,u}$ has exactly the displayed Poisson law, so
$\EE F_{k,r}=q_{k,r}$. The indicators at origins $u,v$ are independent
if $|u-v|\ge h_k$. For the other pairs the upper bound on their covariance is $1/4$:
$\operatorname{Var}(A-B)\ge0$ and
$\operatorname{Var}(A),\operatorname{Var}(B)\le1/4$ give this directly
for indicators $A,B$. Counting at most $N_k(2h_k-1)$ ordered pairs gives
\[
 \Var(F_{k,r})\le\frac{2h_k-1}{4N_k}\ll\frac{h_k}{n_k},
\]
uniformly in $r$. Indeed $h_k/n_k\to0$ gives
$N_k=n_k-h_k+1\ge n_k/2$ eventually, so $\sum_k h_k/N_k<\infty$.
In particular, overlapping windows have not been treated as independent.

Almost surely all runs are finite, so summing
$I_{j,e,\sigma}$ over $(e,\sigma)$ gives $J_{j+1,L_k}$.
The empirical frequency is therefore the same measurable function of
the source field as $F_{k,r}$ is of the target field. For each fixed
$r\in\NN$ and $\eta>0$, \cref{cor:micro-readout} and Chebyshev's
inequality in the target yield
\[
 \PP_f\!\left(\left|\mu_k^f\{r\}-q_{k,r}\right|>\eta\right)
 \le\Delta_k+\frac{2h_k-1}{4N_k\eta^2}.
\]
This is one application to the event defined by the entire empirical
frequency: the error $\Delta_k$ is not multiplied by the number of
origins. The bound is summable. Borel--Cantelli and a countable
intersection over $r\in\NN$ and $\eta=1/m$, $m\ge1$, show that, almost
surely, $\mu_k^f\{r\}-q_{k,r}\to0$ for every $r$.

Finally, $0\le\tau-h_kp_k<p_k\to0$, so
$q_{k,r}\to q_r:=\Pois(\tau)\{r\}$. For any probability $\mu$ on
$\NN$ and any $R\in\NN$, normalization of the masses gives
\[
 \dTV(\mu,q)
 \le\sum_{r=0}^{R}|\mu\{r\}-q_r|+\sum_{r>R}q_r.
\]
First let $k\to\infty$ with $R$ fixed, and then let $R\to\infty$.
This proves the asserted total-variation convergence. Neither an
unbounded source moment nor a joint coupling of the targets at all
scales is used, and Borel--Cantelli requires no independence between
the source events at different scales.
\end{proof}

\paragraph{Scope and comparison.}
The normalization $h_k2^{-L_k}\to\tau$ is essential: the window-size
estimate with an unspecified $O(1)$ alone does not fix the limiting
mean. The inequality $\alpha<1$ leaves room in the arithmetic budget,
whereas $\alpha>0$ makes the target variance summable. This statement
fixes the entire $f$, unlike \cref{cor:quenched}. Najnudel's empirical
block laws \cite{Najnudel2020} use fixed block length.
Pandey--Wang--Xu \cite[arXiv version 4, Theorem~1.2 and footnote~1]{PandeyWangXu2024}
already study a typical multiplicative function from a uniform spatial
origin: their normalized Steinhaus sums have a complex Gaussian limit,
on sets of functions of measure $1-o(1)$ at each size. The present
statement instead counts rare left-maximal binary runs, converges in
total variation to a finite-mean Poisson law, and has a single
full-probability set for each fixed $(\alpha,\tau)$ along the specified
dyadic scales. Poisson genericity samples long words in a fixed sequence
\cite{AlvarezBecherMereb2023,AlvarezEtAlPoissonGenericity}; it is a different empirical
observable. No all-scale or simultaneous all-parameter assertion, and no
assertion about Liouville, follows here. No result about records,
occupation or scale flows is an input. The next observation shows why
the restriction to a count is substantive.

\begin{empiricalobstruction}[A support obstruction for the fixed local field]
\label{rem:empirical-field-obstruction}
Keep the scales of \cref{cor:empirical-poisson}. For any fixed $f$, let
\[
 \mathbf Z_k(f,u)=(J_{u+i+1,L_k}(f))_{1\le i\le h_k},\qquad
 \widehat\mu_k^f=\frac1{N_k}\sum_{u=0}^{N_k-1}\delta_{\mathbf Z_k(f,u)},
 \qquad Q_k=\bigotimes_{i=1}^{h_k}\Pois(p_k).
\]
Then
\[
 \liminf_{k\to\infty}\dTV(\widehat\mu_k^f,Q_k)
       \ge 1-e^{-\tau}(1+\tau)>0.
\]
Indeed, the support $\mathscr S_k(f)$ of $\widehat\mu_k^f$ contains at
most $N_k$ configurations. For large $k$, $p_k\le1$ and every
$z\in\NN^{h_k}$ with $|z|=\sum_i z_i\ge2$ satisfies
\[
 Q_k\{z\}=e^{-h_kp_k}\frac{p_k^{|z|}}{\prod_i z_i!}\le p_k^2.
\]
Testing the complement of $\mathscr S_k(f)$ within $\{|z|\ge2\}$ gives
\[
 \dTV(\widehat\mu_k^f,Q_k)
 \ge 1-e^{-h_kp_k}(1+h_kp_k)-N_kp_k^2.
\]
Now $h_kp_k\to\tau$ and $N_kp_k^2=M_k^{-1+o(1)}\to0$.
This holds for every fixed realization, not merely almost every one.
Thus the count law can converge even though the full local vector,
expressed relative to its sampled origin, cannot have the analogous
product-Poisson approximation in total variation.
\end{empiricalobstruction}
